\section{Free massive fields}\label{sec:free}

Here $d\ge1$, $n\ge2$, $L\ge n+1$, $\mathbb T=(\Z/L\Z)^d$, $\mu>0$ is a lattice
mass, $-\Delta$ is the lattice Laplacian and $Q=-\Delta+\mu^2$. For a positive
definite $Q$ on $\R^N$ let $\gamma_Q$ be the centred Gaussian law with
covariance $C=Q^{-1}$, and for $V\subset\{1,\dots,N\}$ let $\mathcal P_V$ be the
conditional expectation given the coordinates outside $V$. For the cubes
$V_x=x+\{0,\dots,n-1\}^d$, $x\in\mathbb T$, put
\begin{equation}\label{eq:ffvarpi}
\varpi(Q)=\inf_f\frac{n^{-d}\sum_x\|f-\mathcal P_{V_x}f\|^2}{\Var_{\gamma_Q}(f)},
\end{equation}
the infimum over nonconstant $f\in L^2(\gamma_Q)$, the best constant in
$\BP(n,\cdot)$. With $s=n\mu$ and $\mu=ma$, $n=\ell/a$, the parameter $s=m\ell$
is the mass times the physical size of the cube.

\begin{lemma}[linear functions suffice]\label{lem:chaos}
For a positive definite $Q$ and any family of coordinate sets $(V_x)$, the
infimum in~\eqref{eq:ffvarpi} equals its infimum over the linear functions
$\ell_h(\phi)=\ip h\phi$, and
\[
\|\ell_h-\mathcal P_V\ell_h\|^2=\ip{h_V}{Q_{VV}^{-1}h_V},\qquad
\Var(\ell_h)=\ip h{Q^{-1}h},
\]
with $h_V$ the restriction of $h$ to $V$ and $Q_{VV}$ the block of $Q$ on $V$.
\end{lemma}

\begin{proof}
The conditional law of $\phi_V$ given $\phi_{V^c}$ is Gaussian with covariance
$Q_{VV}^{-1}$, which gives the formula for $\ell_h$. Write $\phi=Q^{-1/2}\xi$ with
$\xi$ standard Gaussian. Conditioning on $\phi_{V^c}$ is conditioning on the
Gaussian subspace $K_V=\{\ip u\xi:u\in Q^{-1/2}\R^{V^c}\}$, and on the $m$-th
Wiener chaos, identified with the symmetric tensors $\mathrm{Sym}^m(\R^N)$, it
acts as $\pi_V^{\otimes m}$, with $\pi_V$ the orthogonal projection onto $Q^{-1/2}
\R^{V^c}$~\cite[Ch.~4]{Jan97}. The chaoses are orthogonal and invariant under
every $\mathcal P_V$, so it suffices to treat one chaos $m\ge1$. The projection
$I-\pi_V^{\otimes m}$ dominates $(I-\pi_V)\otimes I^{\otimes(m-1)}$, as both are
projections that commute and the range of the second lies in that of the
first. For a symmetric tensor $F$, therefore,
\[
\sum_x\|F-\pi_{V_x}^{\otimes m}F\|^2\ge\Bigl\langle F,\Bigl(\sum_x(I-\pi_{V_x})\Bigr)
\otimes I^{\otimes(m-1)}F\Bigr\rangle\ge\lambda\|F\|^2,
\]
where $\lambda$ is the least eigenvalue of $\sum_x(I-\pi_{V_x})$ on $\R^N$, which is
the infimum of the quotient over linear functions ($m=1$).
\end{proof}

\begin{theorem}[the exact constant]\label{thm:free}
Let $V=\{0,\dots,n-1\}^d$, $-\Delta_D$ its Dirichlet Laplacian and $\mathbf 1$ its
constant function. Then
\[
\varpi(-\Delta+\mu^2)=\frac{\mu^2}{n^d}\ip{\mathbf 1}{(-\Delta_D+\mu^2)^{-1}\mathbf 1},
\]
attained only at multiples of $\ell_{\mathbf 1}$ among linear functions, and
\[
\Bigl(\frac56\Bigr)^d\frac{s^2}{s^2+10d}\le\varpi(-\Delta+\mu^2)\le\frac{s^2}
{s^2+4dn^2\sin^2\frac\pi{2n+2}}\le\frac{s^2}{s^2+16d/9}.
\]
\end{theorem}

\begin{proof}
Since $L\ge n+1$, the block $Q_{V_xV_x}$ is the translate of $Q_{VV}=-\Delta_D+
\mu^2$; put $G_V=Q_{VV}^{-1}$. By Lemma~\ref{lem:chaos}, $n^d\varpi$ is the least
value of $\ip h{Sh}/\ip h{Q^{-1}h}$ with $S=\sum_xR_{V_x}^{\mathsf T}Q_{V_xV_x}^{-1}
R_{V_x}$, that is, the least eigenvalue of $Q^{1/2}SQ^{1/2}$. Translations commute
with $Q$ and $S$, so the plane waves $e_k(y)=e^{ik\cdot y}$ diagonalise it with
eigenvalues $\lambda_k=\hat q(k)\ip{e_k^V}{G_Ve_k^V}$, where $\hat q(k)=\mu^2+
\sum_i(2-2\cos k_i)$ and $e_k^V$ is the restriction to $V$.

Let $B$ map $f$ to the function $u\mapsto\sum_{w\sim u,\,w\notin V}f(w)$ on $V$.
Then $Q_{VV}R_Vf=R_VQf+Bf$, so $Q_{VV}e_k^V=\hat q(k)e_k^V+Be_k$, and pairing
$G_V$ of this with $e_k^V$ gives
\[
\lambda_k=n^d-\sum_{u,u'\in V}\sum_{w\sim u',\,w\notin V}G_V(u,u')\cos\bigl(k\cdot
(w-u)\bigr).
\]
Write $Q_{VV}=(2d+\mu^2)(I-K)$ with $K$ entrywise nonnegative and $\|K\|\le
2d/(2d+\mu^2)<1$; the Neumann series shows that $G_V$ has positive entries.
Hence $\lambda_k-\lambda_0=\sum G_V(u,u')(1-\cos(k\cdot(w-u)))\ge0$, with equality only
at $k=0$ (take $u=u'$ on a face and $w=u+e_i$). Since $Q\mathbf 1=\mu^2\mathbf 1$,
$\lambda_0=\mu^2\ip{\mathbf 1}{G_V\mathbf 1}$, which is the formula.

\emph{Upper bound.} $G_V\le(\lambda_1+\mu^2)^{-1}$ with $\lambda_1=4d\sin^2\frac\pi
{2n+2}$ the least eigenvalue of $-\Delta_D$, and $\|\mathbf 1\|^2=n^d$; then
$\sin t\ge2t/\pi$ gives $n^2\lambda_1\ge4dn^2/(n+1)^2\ge16d/9$.

\emph{Lower bound.} By the variational formula $\ip{\mathbf 1}{G_V\mathbf 1}\ge
\ip{\mathbf 1}w^2/\ip w{Q_{VV}w}$, take $w(u)=\prod_i\tau(u_i)$ with the parabolic
profile $\tau(s)=6(s+1)(n-s)/(n(n+1)(n+2))$ on $\{0,\dots,n-1\}$, zero
outside. Then $\sum\tau=1$, $S_2=\sum\tau^2=\alpha_n/n$ and $D_2=\sum_{s\in\Z}
(\tau(s+1)-\tau(s))^2=\beta_n/n^3$ with
\[
\alpha_n=\frac{6(n^2+2n+2)}{5(n+1)(n+2)}\le\frac65,\qquad
\beta_n=\frac{12n^2}{(n+1)(n+2)}=10\rho_n\alpha_n,\quad
\rho_n=\frac{n^2}{n^2+2n+2}<1,
\]
by the power sums $\sum_{u=1}^nu(n+1-u)=n(n+1)(n+2)/6$, $\sum_{u=1}^nu^2(n+1-u)^2=
n(n+1)(n+2)(n^2+2n+2)/30$ and $\sum_{u=0}^n(n-2u)^2=n(n+1)(n+2)/3$. The Dirichlet
form of the product is $\ip w{-\Delta_Dw}=dS_2^{d-1}D_2$ and $\ip{\mathbf 1}w=1$, so
\[
\varpi\ge\frac{\mu^2}{n^d(\mu^2S_2^d+dS_2^{d-1}D_2)}=\frac1{\alpha_n^d+d\alpha_n^{d-1}
\beta_ns^{-2}}=\alpha_n^{-d}\frac{s^2}{s^2+10d\rho_n}\ge\Bigl(\frac56\Bigr)^d
\frac{s^2}{s^2+10d}.\qedhere
\]
\end{proof}

At $n=2$ the two bounds coincide with $s^2/(s^2+4d)$. For $d=4$ the lower
bound is $0.0118$ at $s=1$ and $0.0438$ at $s=2$; the upper bound vanishes as
$s\to0$, so a massless free field satisfies no cube condition uniform in the
lattice spacing. The same computation applies, with the Hodge decomposition of the
torus, to the free massive vector field on gauge-invariant functions; it is
not used here.


\subsection{The Casimir defect of the free field}\label{ssec:freecasimir}

The Casimir condition (C) of Proposition~\ref{prop:Zonescale} asks that the
physical side of a spatial torus times its Casimir defect under halving be
at most $\frac1{70}$ at one scale and tend to zero at large scales. For the
free field both quantities can be computed exactly in the continuum limit,
and the computation calibrates the condition: it fails for a massless field
at every scale, and it holds for a field of mass $\mu$ exactly beyond the
scale $17.7/\mu$. Throughout this subsection the spatial torus is
$\Lambda_{\vec\ell}=\prod_{i=1}^3\Z/\ell_i\Z$ with all $\ell_i\ge2$, $V=
\ell_1\ell_2\ell_3$, its dual momenta are $k\in\prod_i\frac{2\pi}{\ell_i}
\Z/\ell_i\Z$, $\hat k^2=\sum_i4\sin^2(k_i/2)$, and $\mu_0\ge0$ is a lattice
mass. The transfer matrix of the free field of mass $\mu_0>0$ is the
operator $\mathbb T_{\vec\ell}$ on $L^2(\R^{\Lambda_{\vec\ell}})$ with the
kernel
\begin{equation}\label{eq:freekernel}
\mathbb T_{\vec\ell}(\varphi,\varphi')=\exp\Bigl(-\tfrac12|\varphi'-\varphi|^2
-\tfrac14\ip\varphi{Q\varphi}-\tfrac14\ip{\varphi'}{Q\varphi'}\Bigr),\qquad
Q=-\Delta+\mu_0^2,
\end{equation}
the spatial part of the action being split symmetrically between the two
slices as in Lemma~\ref{lem:transfer}; $\Tr\mathbb T_{\vec\ell}^M$ is the
partition function $Z(\vec\ell,M)$ of the free field on $\Lambda_{\vec\ell}
\times\Z/M\Z$ for every $M\ge1$, which is symmetric in its four arguments.
Put
\[
\omega(k)=\operatorname{arccosh}\Bigl(1+\frac{\mu_0^2+\hat k^2}2\Bigr),
\qquad E_0(\vec\ell)=-\log\lambda_0(\mathbb T_{\vec\ell}),
\]
and define $\kappa(\vec\ell)=E_0(\vec\ell)-\sum_{s\in\mathcal S}E_0(s)$ and
$\eta(N;s)$ as in Section~\ref{ssec:onescale}, with the eight halves $s\in
\mathcal S$ of sides $\lfloor\ell_i/2\rfloor$ and $\lceil\ell_i/2\rceil$;
there is no coupling and hence no index $p$. Let $p_t(x)=\frac1{2\pi}
\int_{-\pi}^{\pi}e^{iqx-4t\sin^2(q/2)}\,dq=e^{-2t}I_{|x|}(2t)$, $x\in\Z$, be
the heat kernel of the Laplacian of $\Z$, and $h_L(t)=\sum_{j\in\Z/L\Z}
e^{-4t\sin^2(\pi j/L)}=\Tr e^{t\Delta_{\Z/L\Z}}$.

\begin{lemma}[the free vacuum energy]\label{lem:freevacuum}
Let $\mu_0>0$. Then $\mathbb T_{\vec\ell}$ is a positive Hilbert--Schmidt
operator with a continuous strictly positive kernel, its largest eigenvalue
is simple, and
\begin{gather*}
E_0(\vec\ell)=\frac12\sum_k\omega(k)-\frac V2\log2\pi,\\
\kappa(\vec\ell)=\frac12\int_0^\infty\frac{dt}t\,e^{-\mu_0^2t}p_t(0)\Bigl[
\prod_{j=1}^3\bigl(h_{\lfloor\ell_j/2\rfloor}(t)+h_{\lceil\ell_j/2\rceil}(t)
\bigr)-\prod_{j=1}^3h_{\ell_j}(t)\Bigr]\ge0 .
\end{gather*}
The integral converges absolutely also at $\mu_0=0$, and $\kappa_0(\vec\ell)
:=\lim_{\mu_0\downarrow0}\kappa(\vec\ell)$ exists, equals that integral at
$\mu_0=0$, and equals $\frac12\sum_{k\ne0}\omega_0(k)-\frac12\sum_s\sum_{k\ne
0}\omega_0(k)$ with $\omega_0(k)=\operatorname{arccosh}(1+\hat k^2/2)$.
\end{lemma}

\begin{proof}
The kernel~\eqref{eq:freekernel} is $e^{-\ip\varphi{Q\varphi}/4}G(\varphi-
\varphi')e^{-\ip{\varphi'}{Q\varphi'}/4}$ with $G$ a Gaussian, a positive
definite function, so $\mathbb T_{\vec\ell}$ is a positive operator; its
kernel is continuous, strictly positive and square integrable because $Q>0$,
so it is Hilbert--Schmidt and, by Jentzsch's theorem, its largest eigenvalue
is simple with a positive eigenfunction. In the orthonormal basis of
real Fourier modes the quadratic forms $|\varphi'-\varphi|^2$ and
$\ip\varphi{Q\varphi}$ are diagonal, so $\mathbb T_{\vec\ell}$ is the tensor
product over the modes of the one-dimensional operators with kernel
$K_a(x,y)=\exp(-\frac12(x-y)^2-\frac a4(x^2+y^2))$, $a=\mu_0^2+\hat k^2>0$. For
$c>0$ and $b=c-1-a/2$,
\[
\int_\R K_a(x,y)e^{-by^2/2}dy=\sqrt{\frac{2\pi}c}\exp\Bigl(-\frac{x^2}2+
\frac{x^2}{2c}-\frac a4x^2\Bigr),
\]
which equals $\sqrt{2\pi/c}\,e^{-bx^2/2}$ exactly when $c^2-(2+a)c+1=0$, that
is $c=e^{\omega}$ with $\cosh\omega=1+a/2$, and then $b=e^{\omega}-\cosh\omega=\sinh\omega>0$. The positive function $e^{-bx^2/2}$ is therefore an
eigenfunction of the positive kernel $K_a$, hence its Perron eigenfunction,
with eigenvalue $\sqrt{2\pi}\,e^{-\omega/2}$; the largest eigenvalue of the
tensor product is the product, which is the formula for $E_0$.

For the integral representation, $\frac1{2\pi}\int_{-\pi}^\pi\log(A-2\cos q)
\,dq=\log\frac{A+\sqrt{A^2-4}}2$ for $A\ge2$ (write $A-2\cos q=c|1-e^{iq}/c|^2$
with $c+1/c=A$, $c\ge1$, and use $\int_{-\pi}^{\pi}\log|1-e^{iq}/c|\,dq=0$), so with
$A=2+a$,
$\omega(k)=\frac1{2\pi}\int_{-\pi}^{\pi}\log\bigl(a+4\sin^2\tfrac q2\bigr)dq$.
By Frullani's formula $\log\lambda=\int_0^\infty(e^{-t}-e^{-\lambda t})\frac
{dt}t$, applied with $\lambda=a+4\sin^2\frac q2\in[a,a+4]$ and the two
integrals exchanged, which is legitimate because $\lambda$ ranges over a
compact subset of $(0,\infty)$,
\[
\omega(k)=\int_0^\infty\frac{dt}t\Bigl[e^{-t}-e^{-\mu_0^2t}p_t(0)\prod_{j=1}^3
e^{-4t\sin^2(k_j/2)}\Bigr].
\]
Summing over $k$ gives $\sum_k\omega(k)=\int\frac{dt}t[Ve^{-t}-e^{-\mu_0^2t}
p_t(0)\prod_jh_{\ell_j}(t)]$, each bracket being $O(t)$ at $0$ and $O(e^{-t}+
e^{-\mu_0^2t})$ at infinity. The extensive terms cancel in $\kappa$ because
$\sum_sV(s)=V$, and $\sum_s\prod_jh_{s_j}=\prod_j(h_{\lfloor\ell_j/2\rfloor}+
h_{\lceil\ell_j/2\rceil})$, which is the formula. Positivity of $\kappa$ is
the proof of Proposition~\ref{prop:halving}: Lemma~\ref{lem:cutside}, which
uses only the positivity and the Hilbert--Schmidt property of the transfer
matrices in the three spatial directions and the symmetry of $Z$, gives
$Z(\vec\ell,M)\le\prod_sZ(s,M)$, and the $M$-th root gives $\lambda_0(\vec
\ell)\le\prod_s\lambda_0(s)$.

At $\mu_0=0$ the integrand is still absolutely integrable: as $t\to0$, $h_L(t)=
L-2Lt+O(t^2)$ for $L\ge2$ because $\sum_{j}\sin^2(\pi j/L)=L/2$, so the bracket
is $O(t^2)$; as $t\to\infty$ the bracket is bounded by $9\max_j\ell_j^3$ and
$p_t(0)\le(4\pi t)^{-1/2}$ (the bound of Lemma~\ref{lem:freelimit}(ii) below).
Dominated convergence as $\mu_0\downarrow0$ identifies the limit of the
integral with its value at $\mu_0=0$. On the other side, $\omega(0)=
\operatorname{arccosh}(1+\mu_0^2/2)\to0$ and the other $\omega(k)$ are
continuous in $\mu_0$, which gives the mode sum.
\end{proof}

For $v>0$ put $\vartheta(v)=\sum_{j\in\Z}e^{-4\pi^2j^2v}$, $\beta(v)=8
\vartheta(v)^3-\vartheta(v/4)^3$, and for $x\ge0$
\begin{equation}\label{eq:cdef}
c(x)=\frac1{\sqrt{4\pi}}\int_0^\infty v^{-3/2}e^{-x^2v/4}\beta(v)\,dv .
\end{equation}

\begin{proposition}[the continuum Casimir function]\label{prop:cfunction}
\begin{enumerate}[label=\textup{(\roman*)},leftmargin=*]
\item $\beta$ is positive, with the Poisson form
\[
\beta(v)=8(4\pi v)^{-3/2}\sum_{w\in\Z^3\setminus0}\bigl(e^{-|w|^2/4v}-
e^{-|w|^2/v}\bigr),
\]
and $c$ is finite, strictly decreasing on $[0,\infty)$ and tends to $0$ at
infinity.
\item For $x>0$, with $K_2$ the modified Bessel function of the second kind,
\begin{gather*}
c(x)=\frac{x^2}{4\pi^2}\sum_{w\in\Z^3\setminus0}\frac{4K_2(x|w|/2)-K_2(x|w|)}
{|w|^2},\\
c(0)=\frac{15}{2\pi^2}\sum_{w\in\Z^3\setminus0}|w|^{-4}>6.8,\qquad c(0)=12.5630\ldots
\end{gather*}
\item For $x>0$, with $\varepsilon(x)=\sum_{|w|>1}|w|^{-5/2}e^{-x(|w|-1)/2}$,
\[
\frac9{2\pi^{3/2}}\,x^{3/2}e^{-x/2}\Bigl(1+\frac2x\Bigr)\le c(x)\le
\frac6{\pi^{3/2}}\,x^{3/2}e^{-x/2}\Bigl(1+\frac6x+\frac{12}{x^2}\Bigr)\Bigl(1+
\frac{\varepsilon(x)}6\Bigr).
\]
\item The equation $c(x)=\frac1{70}$ has exactly one root $X_*$, and
$16.75\le X_*\le17.99$.
\end{enumerate}
\end{proposition}

\begin{proof}
(i) Poisson summation gives $\vartheta(u)=(4\pi u)^{-1/2}\sum_{j}e^{-j^2/4u}$,
so that $\vartheta(v)^3=(4\pi v)^{-3/2}\sum_{w\in\Z^3}e^{-|w|^2/4v}$ and
$\vartheta(v/4)^3=8(4\pi v)^{-3/2}\sum_{w\in\Z^3}e^{-|w|^2/v}$; the terms
$w=0$ cancel in $\beta$ and every other term is positive. The integrand
of~\eqref{eq:cdef} is positive and strictly decreasing in $x$; it is
integrable because $\beta(v)=O(v^{-3/2}e^{-1/4v})$ at $0$ and $\beta(v)\to7$
at infinity, and dominated convergence gives $c(x)\to0$.

(ii) For $p,q>0$, $\int_0^\infty v^{-3}e^{-pv-q/v}dv=2(p/q)K_2(2\sqrt{pq})$
\cite[\S6.22]{Wat44}, and at $p=0$ the integral is $q^{-2}$. Inserting the
Poisson form of $\beta$ into~\eqref{eq:cdef}, the prefactor is $(4\pi)^{-1/2}
\cdot8(4\pi)^{-3/2}=1/(2\pi^2)$, and the terms $p=x^2/4$, $q=|w|^2/4$ and $q=
|w|^2$ give the series; at $x=0$ each $w$ contributes $(16-1)|w|^{-4}$. The
sum $Z_4=\sum_{w\ne0}|w|^{-4}$ exceeds the six terms $|w|=1$ and the twelve
terms $|w|^2=2$, so $Z_4>9$ and $c(0)>135/(2\pi^2)>6.8$; the value
$Z_4=16.5323\ldots$ gives $c(0)=12.5630\ldots$.\footnote{$Z_4=\int_0^\infty t
\bigl(\theta(t)^3-1\bigr)dt$ with $\theta(t)=\sum_je^{-tj^2}$, by $\int_0^
\infty te^{-t|w|^2}dt=|w|^{-4}$; the integral is computed with the Poisson
form of $\theta$ for $t<1$. Values of $c$ from the Bessel series: $c(8)=
0.787$, $c(12)=0.158$, $c(16)=0.0297$, $c(18)=0.0126$, $c(20)=0.0053$, and
the root of $c=\frac1{70}$ is $X_*=17.7106$. Nothing below rests on these
numbers; the bracket in (iv) is proved from the bounds (iii).}

(iii) From $K_\nu(z)=\int_0^\infty e^{-z\cosh u}\cosh\nu u\,du$, $K_\nu(z)$ is
decreasing in $z>0$ and increasing in $\nu\ge0$, and $K_{3/2}(z)=\sqrt{\pi/2z}
\,e^{-z}(1+1/z)$, $K_{5/2}(z)=\sqrt{\pi/2z}\,e^{-z}(1+3/z+3/z^2)$. For the
lower bound keep the six terms $|w|=1$ of the series, where $4K_2(x/2)-
K_2(x)\ge3K_2(x/2)\ge3K_{3/2}(x/2)$. For the upper bound drop the negative
terms, replace $K_2$ by $K_{5/2}$ and use, for $r=|w|\ge1$ and $z=x/2$,
$K_{5/2}(rz)\le r^{-1/2}e^{-(r-1)z}K_{5/2}(z)$, which is read off from the
closed form; the terms $|w|=1$ give the factor $6$ and the others give
$\varepsilon(x)$.

(iv) Uniqueness is (i). The lower bound of (iii) is $\frac9{2\pi^{3/2}}(x^{3/2}+
2x^{1/2})e^{-x/2}$, decreasing for $x>2$ since $(x^{3/2}+2x^{1/2})'-\frac12
(x^{3/2}+2x^{1/2})=x^{-1/2}(1+\frac x2-\frac{x^2}2)$; its value at $x=16.75$
exceeds $\frac1{70}$ (it is $\frac1{70}$ at $16.7521$), so $c>\frac1{70}$ on
$[0,16.75]$. The upper bound of (iii) is a product of decreasing functions of
$x$ (the factor $(x^{3/2}+6x^{1/2}+12x^{-1/2})e^{-x/2}$ has logarithmic
derivative $-\frac12\,\frac{x^3+3x^2+6x+12}{x^3+6x^2+12x}<0$, and
$\varepsilon$ is decreasing); it equals $\frac1{70}$ at $x=17.9870$, where
$\varepsilon=0.1246$, so $c<\frac1{70}$ on $[17.99,\infty)$.
\end{proof}

\begin{lemma}[two bounds on the heat kernel]\label{lem:freelimit}
\begin{enumerate}[label=\textup{(\roman*)},leftmargin=*]
\item For $0\le t\le2$ and $x\in\Z$, $p_t(x)\le t^{|x|}/|x|!$, and $p_t(0)\le1$
for all $t$.
\item For $t>0$, $p_t(0)\le(\pi/16t)^{1/2}$ and $p_t(x)\le(\pi/16t)^{1/2}
\exp\bigl(-x^2/(4t+2|x|)\bigr)$.
\item $\sqrt{4\pi t}\,p_t(0)\to1$ as $t\to\infty$.
\end{enumerate}
\end{lemma}

\begin{proof}
(i) $p_t(x)=e^{-2t}\sum_{j\ge0}t^{2j+|x|}/(j!(j+|x|)!)\le e^{-2t}e^{t^2}t^{|x|}/
|x|!$, and $t^2\le2t$ for $t\le2$; $p_t(0)\le1$ because $I_0(z)\le e^z$.
(ii) The integrand of $p_t(x)$ is entire and $2\pi$-periodic in $q$, so the
contour may be shifted to $q+i\eta$, $\eta\ge0$, which gives for $x\ge0$
$p_t(x)\le e^{-\eta x}\frac1{2\pi}\int e^{-2t(1-\cos q\cosh\eta)}dq=e^{-\eta x}
e^{2t(\cosh\eta-1)}\,e^{-z}I_0(z)$ with $z=2t\cosh\eta$. Since $\sin(\theta/2)
\ge\theta/\pi$ on $[0,\pi]$,
\[
e^{-z}I_0(z)=\frac1\pi\int_0^\pi e^{-2z\sin^2(\theta/2)}d\theta\le\frac1\pi
\int_0^\infty e^{-2z\theta^2/\pi^2}d\theta=\sqrt{\frac\pi{8z}}\le\sqrt{\frac
\pi{16t}} .
\]
With $\sinh\eta=y:=x/2t$ the exponent is $-2t\Psi(y)$, $\Psi(y)=y\operatorname{
arcsinh}y-\sqrt{1+y^2}+1$, and $\Psi(y)\ge y^2/(2(1+y))$ because both vanish at
$0$ and $\Psi'(y)=\operatorname{arcsinh}y\ge\log(1+y)\ge\frac12(1-(1+y)^{-2})=
\frac d{dy}\frac{y^2}{2(1+y)}$; and $2t\cdot\frac{y^2}{2(1+y)}=\frac{x^2}{4t+2x}$.
(iii) is $I_0(z)\sim e^z/\sqrt{2\pi z}$ \cite[\S7.23]{Wat44}.
\end{proof}

\begin{theorem}[the exact Casimir limit of the free field]\label{thm:freelimit}
\begin{enumerate}[label=\textup{(\roman*)},leftmargin=*]
\item Let $m\to\infty$, $\ell_i\in\{m,m+1\}$ and $\mu_0m\to x\in[0,\infty]$.
Then $(m+1)\kappa(\vec\ell)\to c(x)$, with $c(\infty)=0$.
\item For the free field of physical mass $\mu\ge0$, lattice mass $\mu_0=
\mu/N$, and every $s>0$: $\lim_{N\to\infty}\eta(N;s)=c(\mu s)$.
\end{enumerate}
\end{theorem}

\begin{proof}
(i) Substitute $t=m^2\tau$ in Lemma~\ref{lem:freevacuum}:
\begin{gather*}
m\kappa(\vec\ell)=\frac12\int_0^\infty\frac{d\tau}\tau\,e^{-\mu_0^2m^2\tau}
\bigl(m\,p_{m^2\tau}(0)\bigr)\,\mathcal B_m(\tau),\\
\mathcal B_m(\tau)=
\prod_j\bigl(h_{s_j'}+h_{s_j''}\bigr)(m^2\tau)-\prod_jh_{\ell_j}(m^2\tau),
\end{gather*}
with $s_j'=\lfloor\ell_j/2\rfloor$, $s_j''=\lceil\ell_j/2\rceil$. For fixed
$\tau>0$, $m\,p_{m^2\tau}(0)\to(4\pi\tau)^{-1/2}$ by Lemma~\ref{lem:freelimit}
(iii). If $L/m\to\lambda>0$, then, taking the representatives $|j|\le L/2$,
$h_L(m^2\tau)=\sum_{|j|\le L/2}e^{-4m^2\tau\sin^2(\pi j/L)}\to\sum_je^{-4\pi^2
j^2\tau/\lambda^2}=\vartheta(\tau/\lambda^2)$, by dominated convergence in $j$
with the dominating sequence $e^{-16j^2m^2\tau/L^2}\le e^{-4j^2\tau}$, from
$\sin(\pi j/L)\ge2|j|/L$ and $m/L\ge\frac12$. Here $\lambda=1$ for the sides
and $\lambda=\frac12$ for the halves, so $\mathcal B_m(\tau)\to8\vartheta(4\tau)
^3-\vartheta(\tau)^3$, and $e^{-\mu_0^2m^2\tau}\to e^{-x^2\tau}$ (which is $0$
for $x=\infty$). The substitution $\tau=v/4$ turns the limiting integral into
$c(x)$. It remains to dominate the integrand by a function of $\tau$ alone.
By the Poisson identity $h_L(t)=L\sum_{w\in\Z}p_t(Lw)$, the discrete Fourier
inversion of $\sum_wp_t(x+Lw)=\frac1L\sum_ke^{ikx-4t\sin^2(k/2)}$ at $x=0$,
the terms $w=0$ cancel in $\mathcal B_m$ and
\[
|\mathcal B_m(\tau)|\le\sum_{s\in\mathcal S}V(s)\sum_{w\ne0}\prod_jp_t(s_jw_j)+V
\sum_{w\ne0}\prod_jp_t(\ell_jw_j),\qquad t=m^2\tau,
\]
where every side involved lies in $[m/3,m+1]$. Write $r=\lceil m/3\rceil$ for
the least of them. \emph{Region $t\le1$.} By Lemma~\ref{lem:freelimit}(i),
for a side $L\ge r$ and $k\ge1$, $p_t(Lk)\le t^{Lk}/(Lk)!\le(t^{L}/L!)^k\le\delta
^k$ with $\delta=t^r/r!\le\frac12$, so $\sum_{k\ge1}p_t(Lk)\le2\delta$, and
$\sum_{w\ne0}\prod_jp_t(L_jw_j)\le(p_t(0)+4\delta)^3-p_t(0)^3\le108\,\delta$.
The region contributes at most $m\cdot9(m+1)^3\cdot108\int_0^1t^r\frac{dt}{t\,
r!}=\frac{972\,m(m+1)^3}{r\cdot r!}\to0$. \emph{Region $1\le t\le m^2$, that
is $m^{-2}\le\tau\le1$.} Put $A=(\pi/16\tau)^{1/2}$, so that $m\,p_t(0)\le A$
and $p_t(0)\le A/m$. For a side $L\ge m/3$ and $k\ge1$, Lemma~\ref
{lem:freelimit}(ii) and $a/(b+c)\ge\frac12\min(a/b,a/c)$ give
\[
\frac{L^2k^2}{4t+2Lk}\ge\frac12\min\Bigl(\frac{k^2}{36\tau},\frac{mk}6\Bigr)
\ge\frac k{72\sqrt\tau},
\]
using $m\ge\tau^{-1/2}$ and $\tau\le1$, so $\sum_{k\ge1}p_t(Lk)\le\frac Am
\sum_{k\ge1}e^{-k/72\sqrt\tau}\le\frac Am\,B(\tau)$ with $B(\tau)=\frac{e^{-1/72
\sqrt\tau}}{1-e^{-1/72}}$, and $\sum_{w\ne0}\prod_jp_t(L_jw_j)\le\bigl(\tfrac Am
\bigr)^3\bigl[(1+2B)^3-1\bigr]\le\bigl(\tfrac Am\bigr)^3\,6B(1+2B)^2$. The
integrand is therefore at most $\frac12\cdot A\cdot9\cdot8m^3\cdot(A/m)^3\cdot
6B(1+2B)^2=216\,A^4B(\tau)(1+2B(\tau))^2$, a function of $\tau$ alone which is
integrable against $d\tau/\tau$ on $(0,1]$, since $A^4=(\pi/16\tau)^2$ and
$B(\tau)\le73\,e^{-1/72\sqrt\tau}$. \emph{Region $\tau\ge1$.} For a side
$L\le m+1\le2m$, $h_L(m^2\tau)\le1+2\sum_{j\ge1}e^{-4j^2\tau}\le1+\frac{2e^{-4}}
{1-e^{-4}}\le1.04$, so $|\mathcal B_m|\le9\cdot1.04^3\le10.2$ and the integrand
is at most $5.1\,A\le5.1(\pi/16)^{1/2}\tau^{-1/2}$, integrable against
$d\tau/\tau$ on $[1,\infty)$. Dominated convergence proves (i), as $(m+1)/m
\to1$.

(ii) Lower bound: the supremum defining $\eta(N;s)$ contains the cube of side
$m=2\lceil Ns/2\rceil$, whose halves are cubes, and $\mu_0m=\mu m/N\to\mu s$
as $N\to\infty$, so $\liminf_N\eta(N;s)\ge\lim(m+1)\kappa=c(\mu s)$ by (i). Upper bound: if $\limsup_N\eta(N;s)>c(\mu s)+\epsilon$,
there are $N_k\to\infty$, $m_k\ge N_ks$ and admissible $\vec\ell_k$ with
$(m_k+1)\kappa(\vec\ell_k)\ge c(\mu s)+\epsilon$; along a subsequence $\mu m_k/
N_k\to x\in[\mu s,\infty]$, and (i) gives the limit $c(x)\le c(\mu s)$ because
$c$ is decreasing, a contradiction.
\end{proof}

\begin{corollary}[the Casimir condition for the free field]\label{cor:freeC}
For the free field of physical mass $\mu\ge0$:
\begin{enumerate}[label=\textup{(\roman*)},leftmargin=*]
\item if $\mu=0$, then $\eta(N;s)\ge c(0)>6.8$ (numerically $12.56$) for every
$N$ and every $s$,
so both clauses of \textup{(C)} fail at every scale $s_0$;
\item if $\mu>0$, then \textup{(C)} holds at the scale $s_0$ when $\mu s_0>
X_*$ and fails when $\mu s_0<X_*$; in particular \textup{(C)} requires $\mu s_0
\ge16.75$, that is, a torus of at least $16.75$ Compton wavelengths;
\item the rate $e^{-m_*s/2}$ of Proposition~\ref{prop:FSC} is attained: for
the free field $\lim_N\eta(N;s)=c(\mu s)\ge\frac9{2\pi^{3/2}}(\mu s)^{3/2}
e^{-\mu s/2}$.
\end{enumerate}
\end{corollary}

\begin{proof}
(i) The supremum defining $\eta(N;s)$ contains the even cubes of every side
$m\ge Ns$, along which $(m+1)\kappa_0\to c(0)$ by Theorem~\ref{thm:freelimit}
(i) at $x=0$. (ii) By Theorem~\ref{thm:freelimit}(ii) and the monotonicity
of $c$, $\lim_N\eta(N;s_0)=c(\mu s_0)$ is below $\frac1{70}$ exactly when
$\mu s_0>X_*$, and $\lim_s\lim_N\eta(N;s)=\lim_sc(\mu s)=0$. (iii) is
Proposition~\ref{prop:cfunction}(iii).
\end{proof}

The free gluons of the weak-coupling limit give the same function with a
multiplicity. Consider $\dim G$ independent copies of the Gaussian lattice
Maxwell field: real link variables $A_{x,i}$ on the spatial links, the
transfer kernel $\exp(-\frac12\sum_{x,i}(A'_{x,i}-A_{x,i})^2-\frac14S(A)-
\frac14S(A'))$ with $S(A)=\sum_{x,i<j}F_{ij}(x)^2$, $F_{ij}=\partial_iA_j-
\partial_jA_i$ the lattice curl, acting on functions invariant under the gauge
transformations $A_{x,i}\mapsto A_{x,i}+\chi_{x+e_i}-\chi_x$; the three
constant modes $\sum_xA_{x,i}$ of each colour, which the curl does not see, are
taken on the circle $\R/2\pi\Z$, as they are for the gauge field.

\begin{proposition}[free gluons]\label{prop:freegluons}
In this model $E_0(\vec\ell)=\dim G\bigl[\sum_{k\ne0}\omega_0(k)-\frac{3V}2
\log2\pi\bigr]$, hence $\kappa^{\mathrm{Max}}(\vec\ell)=2\dim G\cdot\kappa_0(
\vec\ell)$ and
\[
\lim_{N\to\infty}\eta^{\mathrm{Max}}(N;s)=2\dim G\cdot c(0)=
\begin{cases}75.38\ldots&G=\SU(2),\\201.00\ldots&G=\SU(3),\end{cases}
\]
at every scale $s$. The free gluons violate \textup{(C)} at every scale.
\end{proposition}

\begin{proof}
In momentum space $F_{ij}(k)=q_iA_j(k)-q_jA_i(k)$ with $q_i=e^{ik_i}-1$, $|q_i|^2
=4\sin^2(k_i/2)$, and Lagrange's identity gives $\sum_{i<j}|q_iA_j-q_jA_i|^2=
|q|^2|A|^2-|\ip qA|^2$. So for $k\ne0$ the form $S$ has the eigenvalue
$\hat k^2$ on the two-dimensional orthogonal complement of $q$ and $0$ on $A
\propto q$, which is the Fourier transform of a gradient, that is, a gauge
direction; the Fourier change of variables is orthogonal. A gauge-invariant
function depends only on the transverse coordinates and the three constant
modes, and the kinetic form $\frac12\sum(A'-A)^2$ is the sum of the kinetic
forms of these orthogonal coordinates, so the transfer operator is the tensor
product of the one-dimensional operators $K_{\hat k^2}$ of Lemma~\ref
{lem:freevacuum} over the two transverse modes of each $k\ne0$, with largest
eigenvalue $\sqrt{2\pi}e^{-\omega_0(k)/2}$ each, and of the three wrapped
Gaussian kernels on the circle, whose Perron eigenfunction is the constant with
eigenvalue $\sqrt{2\pi}$, the mass of the wrapped normalised Gaussian. The
constants are proportional to $V$ and cancel in $\kappa$; the two transverse
modes give $2\kappa_0$ per colour; and Theorem~\ref{thm:freelimit} at $\mu=0$
gives the limit.
\end{proof}

The pressure bound~\eqref{eq:pressureX} of Section~\ref{ssec:pressure}
is also explicit for the free field. It needs the whole spectrum of the
one-mode kernel, not only its Perron eigenvalue.

\begin{lemma}[the spectrum of the one-mode kernel]\label{lem:freemodes}
Let $a>0$, $\cosh\omega=1+a/2$ with $\omega>0$, and $K_a(x,y)=\exp(-\frac12
(x-y)^2-\frac a4(x^2+y^2))$ on $L^2(\R)$. The eigenvalues of $K_a$ are
$\sqrt{2\pi}\,e^{-\omega/2}e^{-j\omega}$, $j=0,1,2,\dots$, each simple, with
eigenfunctions $q_j(x)e^{-bx^2/2}$, $q_j$ a polynomial of degree $j$ and
$b=\sinh\omega$. Hence $K_a$ is trace class and $\Tr K_a^M=(2\pi)^{M/2}e^{-M
\omega/2}/(1-e^{-M\omega})$ for $M\ge1$, and for the transfer
matrix~\eqref{eq:freekernel} of the free field of mass $\mu_0>0$ on
$\Lambda_{\vec\ell}$,
\[
\log\mathcal Z_{\vec\ell}(M)=-\sum_k\log\bigl(1-e^{-M\omega(k)}\bigr),\qquad
M\ge1 ,
\]
and every eigenvalue $\lambda\ne\lambda_0$ of $\mathbb T_{\vec\ell}$
satisfies $\lambda\le e^{-\omega(0)}\lambda_0$, with equality for the one-mode
excitation of the constant mode.
\end{lemma}

\begin{proof}
Put $c=e^\omega$, so that $b=c-1-a/2$ as in Lemma~\ref{lem:freevacuum}. For a
polynomial $q$, completing the square in $y$,
\[
\int_\R K_a(x,y)\,q(y)e^{-by^2/2}dy=e^{-bx^2/2}\int_\R e^{-\frac c2(y-x/c)^2}
q(y)\,dy=\sqrt{\frac{2\pi}c}\,e^{-bx^2/2}\,\mathbb E\bigl[q(x/c+\xi/\sqrt c)
\bigr],
\]
$\xi$ a standard Gaussian variable, because $-\frac12(x-y)^2-\frac a4(x^2+y^2)
-\frac b2y^2=-\frac c2(y-x/c)^2-\frac b2x^2$ when $c=1+a/2+b$ and $c^2-(2+a)c
+1=0$, which is the case for $c=e^\omega$, $b=\sinh\omega$. For $q(y)=y^j$
the right side is $\sqrt{2\pi/c}\,e^{-bx^2/2}$ times a polynomial of degree
$j$ with leading coefficient $c^{-j}$. So the finite-dimensional space of
polynomials of degree at most $j$ times $e^{-bx^2/2}$ is invariant, and $K_a$
acts on it by a triangular matrix with the diagonal entries $\sqrt{2\pi/c}\,
c^{-i}=\sqrt{2\pi}e^{-\omega/2}e^{-i\omega}$, $0\le i\le j$, which are distinct;
hence there are eigenfunctions $q_je^{-bx^2/2}$ with $\deg q_j=j$ and these
eigenvalues. Since $K_a$ is self-adjoint, eigenfunctions of distinct
eigenvalues are orthogonal, and the $q_je^{-bx^2/2}$ span the polynomials
times $e^{-bx^2/2}$, which are dense in $L^2(\R)$; so they form an orthogonal
basis of eigenfunctions, every eigenvalue is simple, and there are no others.
The trace is the geometric series. The transfer matrix is the tensor product
over the $V$ real Fourier modes of the $K_{a(k)}$, $a(k)=\mu_0^2+\hat k^2$, so
its eigenvalues are $\prod_k\sqrt{2\pi}e^{-\omega(k)/2}e^{-j_k\omega(k)}$ over
all $(j_k)\in\N_0^V$, the largest being $j\equiv0$ and the second $j=1$ at the
$k$ with the least $\omega(k)$, which is $k=0$; dividing by $\lambda_0^M$ and
taking the trace gives the mode sum.
\end{proof}

For $\vartheta>0$ and $x>0$ put
\begin{equation}\label{eq:Fdef}
F_\vartheta(x)=\sum_{n\in\Z^3}-\log\Bigl(1-\exp\bigl(-\vartheta\sqrt{x^2+4\pi^2
|n|^2}\bigr)\Bigr),\qquad F=F_{1/2},
\end{equation}
a convergent series of positive terms, strictly decreasing in $x$ and in
$\vartheta$.

\begin{theorem}[the pressure bound of the free field]\label{thm:pressfree}
Let $\mu>0$ and $s>0$, put $x=\mu s$, and for $N$ with $c=Ns$ an even integer
let the free field of mass $\mu_0=\mu/N$ live on the cube of side $c$; let
$X_N=\log\mathcal Z_{(c,c,c)}(c/2)$ and $D_N=D_{(c,c,c)}(c)$.
\begin{enumerate}[label=\textup{(\roman*)},leftmargin=*]
\item $X_N=\sum_{n\in(-c/2,c/2]^3}-\log\bigl(1-e^{-\frac c2\omega(2\pi n/c)}
\bigr)$, and $\log(1+D_N)$ is the same sum with $\frac c2$ replaced by $c$.
\item If $N\ge\mu/2$, then $X_N\le F_{0.2297}(x)$ and $\log(1+D_N)\le
F_{0.4595}(x)$.
\item $\lim_{N\to\infty}X_N=F(x)$ and $\lim_{N\to\infty}D_N=e^{F_1(x)}-1$; and
$F(x)=\log(2/x)+O(1)$ as $x\downarrow0$.
\item On every spatial torus the second eigenvalue of the time transfer
matrix is $e^{-\Delta_N/N}\lambda_0$ with $\Delta_N=2N\operatorname{arsinh}
(\mu/2N)$, which increases to $\mu$.
\end{enumerate}
Consequently the pressure bound~\eqref{eq:pressureX} holds at every scale and
every mass, with $c_b=\sup_NX_N<\infty$; the clause $\Delta_Ns\ge2(X_N+\log
10)$ of Corollary~\ref{cor:gapbridge}\textup{(ii)}, with the actual gap in
place of $\Delta_E$, holds for all large $N$ if $x>x_*$ and fails for all
large $N$ if $x<x_*$, where $x_*=5.2063\ldots$ is the root of $x/2=F(x)+\log
10$; and \textup{(Z$_0^{\rm dy}$)}, $D_N\le\frac1{10}$, holds for all large
$N$ if $x>x_1$ and fails for all large $N$ if $x<x_1$, where $x_1=2.4851\ldots$
is the root of $F_1(x)=\log1.1$.
\end{theorem}

\begin{proof}
(i) is Lemma~\ref{lem:freemodes} with $M=c/2$ and $M=c$, the dual momenta of
the cube being $2\pi n/c$ with $n\in(-c/2,c/2]^3$.

(ii) For $|k_i|\le\pi$, $\sin^2(k_i/2)\ge k_i^2/\pi^2$, so $\mu_0^2+\hat k^2\ge
(4/\pi^2)(\mu_0^2+|k|^2)$. Put $u=\frac12\sqrt{\mu_0^2+\hat k^2}$; as $\mu_0\le
2$ when $N\ge\mu/2$, $u\le\frac12\sqrt{4+12}=2$, and the concavity of
$\operatorname{arsinh}$ on $[0,\infty)$ with $\operatorname{arsinh}0=0$ gives
$\operatorname{arsinh}u\ge\frac u2\operatorname{arsinh}2$. Since $\omega=
\operatorname{arccosh}(1+2u^2)=2\operatorname{arsinh}u$,
\[
\omega(k)\ge\operatorname{arsinh}(2)\,u\ge\frac{\operatorname{arsinh}2}{\pi}
\sqrt{\mu_0^2+|k|^2}\ge0.4595\sqrt{\mu_0^2+|k|^2},
\]
and with $c\mu_0=x$, $c|k|=2\pi|n|$ this gives $\frac c2\omega\ge0.2297\sqrt{
x^2+4\pi^2|n|^2}$ and $c\,\omega\ge0.4595\sqrt{x^2+4\pi^2|n|^2}$. As $y\mapsto
-\log(1-e^{-y})$ is decreasing and the sums in (i) run over part of $\Z^3$,
(ii) follows.

(iii) For fixed $n$, $c^2\bigl(\mu_0^2+\hat k^2\bigr)\to x^2+4\pi^2|n|^2$ with
$k=2\pi n/c$, and $\omega=2\operatorname{arsinh}(\frac12\sqrt a)=\sqrt a\,(1+
O(a))$, so $\frac c2\omega\to\frac12\sqrt{x^2+4\pi^2|n|^2}$ and $c\,\omega\to
\sqrt{x^2+4\pi^2|n|^2}$. Dominated convergence, with the termwise bounds of
(ii), gives the limits. The term $n=0$ of $F$ is $-\log(1-e^{-x/2})=\log(2/x)+
O(x)$, and the others are bounded uniformly for $x\le1$.

(iv) By Lemma~\ref{lem:freemodes} the second eigenvalue is $e^{-\omega(0)}
\lambda_0$ with $\omega(0)=2\operatorname{arsinh}(\mu_0/2)$, so $\Delta_N=N
\omega(0)$; and $y^{-1}\operatorname{arsinh}y$ decreases on $(0,\infty)$ to
the limit $1$ at $0$, so $\Delta_N=\mu\,y^{-1}\operatorname{arsinh}y$ with $y=
\mu/2N$ increases to $\mu$.

For the consequences: by (ii), $\sup_{N\ge\mu/2}X_N\le F_{0.2297}(x)$, and the
finitely many smaller $N$ give finite values, so $c_b<\infty$. Next, $\Delta_N
s-2(X_N+\log10)\to x-2F(x)-2\log10$, which is strictly increasing in $x$, since
$F$ decreases, tends to $-\infty$ at $0$ and to $+\infty$ at $\infty$; so it
has one root $x_*$ and the sign of the limit decides for all large $N$. In
the same way $D_N\to e^{F_1(x)}-1$ is strictly decreasing in $x$, from
$+\infty$ to $0$, and crosses $\frac1{10}$ once, at $x_1$.
\end{proof}

The same three thresholds, $17.71$ for (C) by Corollary~\ref{cor:freeC},
$5.21$ for the clause of Corollary~\ref{cor:gapbridge}(ii) and $2.49$ for
(Z$_0^{\rm dy}$) itself, measure the slack of the three routes on the one
model where everything is explicit: the Casimir route pays for its
generality in the scale it needs, the dyadic route with a pressure bound
pays a factor of two, and (Z$_0^{\rm dy}$) is what the free field actually
provides beyond two and a half correlation lengths.

For the gauge theory itself nothing of this is proved. At weak coupling the
analysis of L\"uscher~\cite{Lus83} (see also~\cite{KvB87}) writes the vacuum
energy of a torus of side $\ell$, for $\ell\Lambda\ll1$, as the sum of the free
gluon energies and a zero-mode energy $\varepsilon_0g^{2/3}/\ell+O(g^{4/3}/
\ell)$ with $\varepsilon_0>0$, the ground energy of the quartic Hamiltonian of
the constant modes; the latter enters $(m+1)\kappa$ with the coefficient $1-
8\cdot2=-15$, so that in this regime $\eta$ is $2\dim G\,c(0)$ decreased by a
term of order $g^{2/3}$, far above $\frac1{70}$. We do not control the
exchange of the weak-coupling and large-volume limits on the lattice, and
the statement is recorded as a heuristic only. Its consequence is a
delimitation. For a theory whose lightest states have masses $m_i$ and
multiplicities $n_i$, the free-field picture gives $\eta\approx\sum_in_i\,c(m_i
s)$, so that (C) at the scale $s_0$ needs $m_{\mathrm{gap}}s_0\gtrsim17.7$: the
Casimir condition is a statement about tori of about eighteen correlation
lengths, beyond perturbation theory and of the same strength as a finite-size
bound on the gap itself. This is consistent with Theorem~\ref{thm:gapcasimir},
where the condition is derived from the gap on the tori of the halving chain,
and with Corollary~\ref{cor:freeC}(iii), where its rate is that of the gap.
